\chapter{Literature Review}

\section{Introduction}

This chapter reviews the theoretical foundations and recent research related to Zero Trust Architecture (ZTA), Internet of Things (IoT) security, trust management systems, cross-domain authentication and authorization, and Federated Learning (FL). The chapter critically analyzes recent studies and identifies existing limitations that motivate the development of the proposed Federated Learning-Based Adaptive Trust Management Framework for Cross-Domain Zero Trust IoT Networks.


\section{Zero Trust Architecture (ZTA)}

\subsection{Evolution of Network Security Models}

Traditional Perimeter-Based Security Models (PBSM) assume that entities inside a network boundary can be trusted. However, cloud computing, remote work, mobile devices, and IoT ecosystems have rendered perimeter-based defenses insufficient.


The concept of Zero Trust was first introduced by John Kindervag in 2010 and later standardized by NIST SP 800-207. The core principle of ZTA is “Never Trust, Always Verify.”


\paragraph{Core Principles of ZTA}

Continuous Authentication


Continuous Authorization


Least Privilege Access


Micro-Segmentation


Continuous Monitoring


Dynamic Trust Evaluation


\subsection{Benefits of ZTA in IoT Environments}

ZTA addresses major IoT security challenges including:


Insider threats


Lateral movement attacks


Device impersonation


Credential theft


Unauthorized access


However, implementing ZTA in large-scale IoT environments introduces significant challenges related to scalability, trust management, and authentication overhead.


\section{Trust Management in Zero Trust Environments}

\subsection{Concept of Trust}

Trust represents the degree of confidence that a device, user, or service will behave according to security expectations.


Trust evaluation commonly utilizes:


\paragraph{Behavioral Trust}

\paragraph{Historical Trust}

\paragraph{Contextual Trust}

Reputation-Based Trust


\subsection{Existing Trust Models}

Existing trust models can be classified into:


\paragraph{Centralized Trust Models}

Advantages:


Simplicity


Disadvantages:


Single Point of Failure


Limited Scalability


\paragraph{Distributed Trust Models}

Advantages:


Scalability


Fault Tolerance


Disadvantages:


Complex Trust Coordination


\subsection{Challenges of Current Trust Models}

Current trust models suffer from:


Static trust calculation


Lack of trust migration


Poor cross-domain interoperability


Privacy concerns


High computational complexity


\section{Cross-Domain Authentication and Authorization}

\subsection{Cross-Domain Security Requirements}

Modern IoT systems frequently operate across:


Enterprise Networks


Cloud Environments


Industrial Control Systems


Service Provider Networks


Cross-domain interaction requires:


Identity Verification


Trust Transfer


Authorization Coordination


\subsection{Existing Cross-Domain Authentication Approaches}

\paragraph{One-Time Authentication}

Advantages:


Low authentication overhead


Limitations:


Session hijacking


Lateral movement attacks


\paragraph{Blockchain-Based Authentication}

Advantages:


Decentralization


Transparency


Limitations:


High latency


High computational overhead


Energy consumption


\paragraph{Zero Trust Authentication}

Advantages:


Continuous verification


Dynamic authorization


Limitations:


Repeated authentication costs


Trust portability limitations


\section{Federated Learning in Cybersecurity}

\subsection{Overview of Federated Learning}

Federated Learning is a decentralized machine learning paradigm where multiple participants collaboratively train a model without sharing raw data.


Benefits include:


Privacy Preservation


Reduced Data Exposure


Distributed Intelligence


\subsection{Federated Learning for Security Applications}

FL has been applied in:


Intrusion Detection Systems


Malware Detection


Anomaly Detection


Trust Evaluation


\subsection{Challenges of FL}

Model Poisoning Attacks


Data Poisoning


Byzantine Nodes


Non-IID Data Distribution


Aggregation Complexity


\section{Review of Recent Zero Trust Research}

\subsection{Zero Trust Architecture: A Systematic Literature Review (Gambo \& Almulhem, 2025)}

\paragraph{Main Contributions}

Comprehensive taxonomy of ZTA applications.


Identification of enabling technologies.


Analysis of adoption challenges.


\paragraph{Key Findings}

Strong growth of ZTA adoption.


Significant challenges in trust management.


Limited cross-domain implementations.


\paragraph{Identified Gaps}

Lack of trust management frameworks.


Privacy-preserving trust evaluation remains underexplored.


Scalability issues in large-scale IoT systems.


\subsection{Zero Trust Architecture in the Modern Threat Landscape (2026)}

\paragraph{Main Contributions}

Analysis of vulnerabilities and attack surfaces.


Classification of Zero Trust weaknesses.


\paragraph{Key Findings}

Identity systems remain primary attack targets.


Policy engines become critical points of failure.


Cross-domain orchestration remains unsolved.


\paragraph{Identified Gaps}

Absence of unified trust orchestration.


Lack of standard evaluation frameworks.


Limited real-world deployments.


\subsection{DAGA: Dynamic Authentication and Granularized Authorization (2026)}

\paragraph{Main Contributions}

Cross-domain pre-authentication.


LSTM-based trust prediction.


Granularized authorization.


\paragraph{Key Findings}

Reduced authentication latency.


Improved authorization efficiency.


\paragraph{Limitations}

Trust migration remains partially addressed.


Dependence on simulation-based evaluation.


No protection against FL poisoning attacks.


\section{Comparative Analysis of Existing Studies}

\begin{table}[htbp]
\centering
\caption{Comparative Analysis of Existing Studies}
\label{tab:comparative-studies}
\small
\begin{tabularx}{\textwidth}{|>{\raggedright\arraybackslash}p{2.6cm}|>{\raggedright\arraybackslash}p{2.5cm}|>{\raggedright\arraybackslash}p{2.2cm}|X|X|}
\hline
 Study & Main Focus & Technology Used & Strengths & Limitations \\ \hline
 Gambo \& Almulhem (2025) & ZTA Review & SLR & Comprehensive taxonomy & No proposed framework \\ \hline
 Modern Threat Landscape (2026) & Vulnerabilities & SLR & Threat classification & No implementation model \\ \hline
 DFL-ZTA (2025) & Cross-Domain Authentication & DFL & Privacy preservation & Limited trust orchestration \\ \hline
 DAGA (2026) & Adaptive Authentication & FL + LSTM & Reduced latency & Limited trust migration \\ \hline
 Proposed Study & Adaptive Trust Management & FL + Dynamic Trust & Cross-domain trust orchestration & To be evaluated \\ \hline
\end{tabularx}
\end{table}

\section{Research Gap Analysis}

Based on the reviewed literature, several gaps remain unresolved:


Lack of adaptive trust management frameworks capable of operating across multiple security domains.


Absence of effective trust migration mechanisms between heterogeneous IoT domains.


Limited integration between Federated Learning and Zero Trust trust orchestration.


Insufficient privacy-preserving trust evaluation techniques.


Lack of scalable trust coordination mechanisms for large-scale IoT deployments.


Existing solutions focus on authentication optimization rather than holistic trust management.


Therefore, a Federated Learning-Based Adaptive Trust Management Framework is proposed to address these limitations.


\section{Conceptual Framework}

The proposed framework is founded on the integration of:


Zero Trust Architecture


Federated Learning


Dynamic Trust Evaluation


Cross-Domain Trust Orchestration


Adaptive Authorization


These components collectively support secure, scalable, and privacy-preserving trust management across heterogeneous IoT environments.


\section{Chapter Summary}

This chapter reviewed the theoretical foundations and contemporary literature related to Zero Trust Architecture, trust management, cross-domain authentication, and Federated Learning. The review identified several limitations in existing approaches, particularly regarding adaptive trust management and cross-domain trust orchestration. These limitations establish the research gap addressed by the proposed framework. The next chapter presents the research design and system architecture of the proposed solution.
